\documentclass[11pt,a4paper]{article}
\usepackage[utf8]{inputenc}
\usepackage[T1]{fontenc}
\usepackage[spanish]{babel}
\addto\captionsspanish{\renewcommand{\tablename}{Tabla}}
\usepackage{longtable,array}
\newcolumntype{L}[1]{>{\raggedright\arraybackslash}p{#1}}
\usepackage[margin=2.5cm]{geometry}
\usepackage{hyperref}

\begin{document}

\section{Propiedades de calidad y actores del proceso}

\subsection{Marco de evaluación}

La calidad de cada uno de los quince requisitos de la
matriz de identificación de requisitos se evaluó conforme a
ISO/IEC/IEEE 29148 \cite{iso29148}, complementada con
Pohl \cite{pohl2010} y SWEBOK v4 \cite{swebok4}. La norma
define características que debe tener un requisito bien
formulado (necesario, no ambiguo, completo, singular,
factible y verificable) y la consistencia como característica
del conjunto; la trazabilidad se trata como atributo del
requisito. Las ocho propiedades de la práctica se adaptaron
a esas características según se resume en la
Tabla~\ref{tab:propiedades}.

\begingroup
\footnotesize
\setlength{\tabcolsep}{3pt}
\renewcommand{\arraystretch}{1.15}
\begin{longtable}{|L{0.17\textwidth}|L{0.09\textwidth}|L{0.62\textwidth}|}
\caption{Propiedades de calidad y criterio de evaluación aplicado}
\label{tab:propiedades}\\
\hline
\textbf{Propiedad} & \textbf{Abrev.} & \textbf{Criterio aplicado} \\
\hline
\endfirsthead

\multicolumn{3}{c}{\textit{Continuación de la Tabla~\ref{tab:propiedades}}}\\
\hline
\textbf{Propiedad} & \textbf{Abrev.} & \textbf{Criterio aplicado} \\
\hline
\endhead

Necesidad & Nec. &
El requisito define una capacidad, característica o restricción
que responde a una necesidad identificada; no es un añadido
prescindible. \\
\hline

Ausencia de ambigüedad & Amb. &
El enunciado admite una sola interpretación; no usa términos
vagos ni sin definir (``adecuado'', ``aplicable'', ``seleccionados''
con dos sentidos). \\
\hline

Completitud & Comp. &
El enunciado contiene la información necesaria para
implementarlo y comprobarlo (datos, condiciones, límites,
comportamiento esperado). \\
\hline

Consistencia & Cons. &
No contradice a otros requisitos ni usa términos distintos para
el mismo concepto. \\
\hline

Verificabilidad & Ver. &
Existe un procedimiento finito (inspección, prueba o medición)
que permite comprobar si se cumple. \\
\hline

Factibilidad & Fac. &
Es alcanzable con la tecnología, el alcance y la información
disponibles para el equipo. \\
\hline

Trazabilidad & Tra. &
Puede vincularse con una necesidad de origen y con una evidencia
o prueba; tiene identificador único. \\
\hline

Singularidad & Sin. &
Expresa una sola capacidad, característica o restricción; no
agrupa varios requisitos en uno. \\
\hline

\multicolumn{3}{@{}L{0.9\textwidth}@{}}{\footnotesize\textit{Fuente: Elaboración propia a partir de \cite{iso29148,pohl2010,swebok4}.}}\\
\end{longtable}
\endgroup

Cada propiedad se calificó con tres niveles: \textbf{C}
(cumple), \textbf{P} (cumple parcialmente) y \textbf{N}
(no cumple). El resultado global de cada requisito se asignó
con esta regla: \emph{Cumple} si todas las propiedades son C;
\emph{Cumple con ajustes} si tiene una o dos calificaciones P;
y \emph{Requiere corrección} si tiene tres o más P o
cualquier N.

La evaluación se basó en el texto de cada requisito, en las
necesidades identificadas y en las capturas de la exploración
(Figs.~1 a 9). No se midió rendimiento,
disponibilidad ni accesibilidad, por lo que esos aspectos se
evalúan como requisitos escritos, no como funcionamiento del
sistema.

\subsection{Matriz de evaluación de calidad}

La Tabla~\ref{tab:matriz_calidad} resume la calificación de
los quince requisitos y la Tabla~\ref{tab:justificaciones}
presenta la justificación verificable y las correcciones
propuestas para cada uno.

\begingroup
\footnotesize
\setlength{\tabcolsep}{2.5pt}
\renewcommand{\arraystretch}{1.15}

\begin{longtable}{|L{0.09\textwidth}|*{8}{>{\centering\arraybackslash}p{0.062\textwidth}|}L{0.2\textwidth}|}
\caption{Matriz de evaluación de calidad de los requisitos}
\label{tab:matriz_calidad}\\
\hline
\textbf{ID} & \textbf{Nec.} & \textbf{Amb.} & \textbf{Comp.} &
\textbf{Cons.} & \textbf{Ver.} & \textbf{Fac.} & \textbf{Tra.} &
\textbf{Sin.} & \textbf{Resultado} \\
\hline
\endfirsthead

\multicolumn{10}{c}{\textit{Continuación de la Tabla~\ref{tab:matriz_calidad}}}\\
\hline
\textbf{ID} & \textbf{Nec.} & \textbf{Amb.} & \textbf{Comp.} &
\textbf{Cons.} & \textbf{Ver.} & \textbf{Fac.} & \textbf{Tra.} &
\textbf{Sin.} & \textbf{Resultado} \\
\hline
\endhead

RF-01 & C & P & P & C & C & C & C & C & Cumple con ajustes \\ \hline
RF-02 & C & P & P & C & C & C & C & C & Cumple con ajustes \\ \hline
RF-03 & C & C & C & C & C & C & C & C & Cumple \\ \hline
RF-04 & C & P & P & C & C & C & C & C & Cumple con ajustes \\ \hline
RF-05 & C & P & P & C & C & C & C & C & Cumple con ajustes \\ \hline
RF-06 & C & P & P & C & C & C & C & C & Cumple con ajustes \\ \hline
RF-07 & C & P & C & C & C & C & C & C & Cumple con ajustes \\ \hline
RF-08 & C & P & P & C & C & C & C & C & Cumple con ajustes \\ \hline
RNF-01 & C & P & P & C & P & C & P & C & Requiere corrección \\ \hline
RNF-02 & C & P & P & C & P & C & P & C & Requiere corrección \\ \hline
RNF-03 & C & P & P & C & C & C & P & C & Requiere corrección \\ \hline
RNF-04 & C & C & P & C & C & P & P & P & Requiere corrección \\ \hline
RES-01 & C & P & P & C & C & C & C & C & Cumple con ajustes \\ \hline
RES-02 & C & P & P & C & P & C & P & C & Requiere corrección \\ \hline
RES-03 & C & C & C & C & C & C & C & P & Cumple con ajustes \\ \hline
\textbf{Total C} & 15 & 3 & 3 & 15 & 12 & 14 & 10 & 13 & \textbf{1 Cumple} \\ \hline
\textbf{Total P} & 0 & 12 & 12 & 0 & 3 & 1 & 5 & 2 & \textbf{9 con ajustes} \\ \hline
\textbf{Total N} & 0 & 0 & 0 & 0 & 0 & 0 & 0 & 0 & \textbf{5 con corrección} \\ \hline

\multicolumn{10}{@{}L{0.9\textwidth}@{}}{\footnotesize\textit{C = cumple; P = cumple parcialmente; N = no cumple. Fuente: Elaboración propia a partir de la matriz de identificación de requisitos y de las evidencias de la exploración.}}\\
\end{longtable}
\endgroup

Las propiedades con más calificaciones parciales son la
completitud y la ausencia de ambigüedad (12 de 15 cada una),
seguidas de la trazabilidad (5 de 15). No se encontró
ninguna contradicción entre requisitos, por lo que la
consistencia cumple en los quince casos.

\begingroup
\footnotesize
\setlength{\tabcolsep}{2.5pt}
\renewcommand{\arraystretch}{1.15}

\begin{longtable}{|L{0.075\textwidth}|L{0.11\textwidth}|L{0.40\textwidth}|L{0.31\textwidth}|}
\caption{Justificación verificable y correcciones propuestas}
\label{tab:justificaciones}\\
\hline
\textbf{ID} & \textbf{Resultado} & \textbf{Justificación verificable} &
\textbf{Observaciones o corrección necesaria} \\
\hline
\endfirsthead

\multicolumn{4}{c}{\textit{Continuación de la Tabla~\ref{tab:justificaciones}}}\\
\hline
\textbf{ID} & \textbf{Resultado} & \textbf{Justificación verificable} &
\textbf{Observaciones o corrección necesaria} \\
\hline
\endhead

RF-01 & Cumple con ajustes &
Define una sola operación y se vincula con la necesidad
``Localizar productos mediante palabras clave''. Evidencia: búsqueda
de ``Arroz'' (Fig.~2). No indica sobre qué datos
se busca ni qué ocurre cuando no hay coincidencias. &
Reformular: ``El sistema deberá permitir buscar productos por
palabra clave aplicada al nombre y la marca del producto y mostrar
el listado de coincidencias; si no existen coincidencias, deberá
mostrar un mensaje informativo.'' Confirmar en el sitio los
campos que realmente se buscan. \\
\hline

RF-02 & Cumple con ajustes &
Se vincula con la necesidad ``Explorar artículos organizados por
categorías''. Evidencia: Figs.~3 y~4. La
captura de la categoría Supermercado muestra filtros por Categoría,
Precio y Fabricante y un total de 2948 artículos, que el enunciado
no recoge; ``navegar'' no precisa qué se presenta al elegir una
categoría. &
Reformular: ``El sistema deberá permitir seleccionar una categoría
desde el menú de categorías y mostrar los productos que contiene,
con opciones de filtrado por categoría, precio y fabricante.''
Si se quiere conservar la singularidad, separar el filtrado en
un requisito propio. \\
\hline

RF-03 & Cumple &
Enuncia los tres datos a mostrar (nombre, precio e imagen) y se
comprueba por inspección visual en la Fig.~5: nombre
``Arroz Real 2 kg'', imagen del producto y precio de \$3,73.
Se vincula con la necesidad ``Consultar información de los
productos''. &
La verificación se hizo en un solo producto; conviene repetirla
en una muestra de productos de varias categorías. Opcionalmente
incluir la marca, que también aparece en la ficha del producto. \\
\hline

RF-04 & Cumple con ajustes &
Operación única vinculada con ``Agregar productos al carrito de
compras''. Evidencia: mensaje de confirmación de producto agregado
(Fig.~6). La ficha del producto incluye un selector de
cantidad que el enunciado no menciona, ni tampoco la confirmación
al usuario. &
Reformular: ``El sistema deberá permitir agregar al carrito un
producto en la cantidad indicada por el usuario y mostrar un
mensaje de confirmación.'' \\
\hline

RF-05 & Cumple con ajustes &
Define la consulta del carrito, vinculada con la necesidad de
seleccionar productos de interés. Evidencia: Fig.~7
(productos, cantidades y precios). El enunciado no especifica qué
datos se muestran. &
Reformular: ``El sistema deberá mostrar en el carrito, para cada
producto agregado, su nombre, cantidad y precio, y el total de la
compra.'' \\
\hline

RF-06 & Cumple con ajustes &
Operación de actualización vinculada con ``Modificar las
cantidades de productos en el carrito''. Evidencia:
Fig.~8, donde cambian las cantidades y los importes.
El enunciado no indica límites de cantidad ni el recálculo de
importes. &
Reformular: ``El sistema deberá permitir modificar la cantidad
de unidades de un producto en el carrito y recalcular el importe
del producto y el total.'' Definir y verificar los límites de
cantidad (mínimo y máximo). \\
\hline

RF-07 & Cumple con ajustes &
Operación única, verificada en la Fig.~9, donde el
carrito se actualiza tras eliminar un producto. El término
``seleccionados'' se usa en RF-04 para productos elegidos del
catálogo y aquí para productos del carrito, lo que admite dos
lecturas. &
Reformular: ``El sistema deberá permitir eliminar del carrito un
producto agregado previamente.'' Unificar la terminología en
RF-04 a RF-07 (``agregado'' para lo que está en el carrito). \\
\hline

RF-08 & Cumple con ajustes &
Vinculado con ``Consultar la tienda seleccionada''. Evidencia:
Figs.~1 y~9 y la ficha del producto, que
muestra el retiro en Tienda Tía San Carlos. No precisa dónde se
muestra ni qué datos de la tienda se presentan; la ficha del
producto incluye además el enlace ``Cambiar Tienda'', que no está
cubierto. &
Reformular: ``El sistema deberá mostrar el nombre de la tienda
seleccionada por el usuario en la barra superior del sitio.''
Crear un requisito independiente para cambiar de tienda. \\
\hline

RNF-01 & Requiere corrección &
La meta es cuantificable (95\,\% de las búsquedas en un máximo
de 3 s), pero depende de ``condiciones de prueba definidas'' que
no se definen en el documento (red, dispositivo, carga, palabras
clave): dos evaluadores obtendrían resultados no comparables.
Tampoco se vincula a una de las nueve necesidades ni a RF-01
por identificador. Estado en la matriz: ``Propuesto; no medido''. &
Definir las condiciones, por ejemplo: ``Con una conexión de
velocidad definida, un navegador de escritorio y un conjunto
fijo de búsquedas de prueba, el 95\,\% de las búsquedas de RF-01
deberá mostrar resultados en un máximo de 3 s.'' Registrar la
necesidad de origen. \\
\hline

RNF-02 & Requiere corrección &
Una disponibilidad mensual de 99,5\,\% equivale a unas
3,6 horas de indisponibilidad en un mes de 30 días
($0{,}5\,\% \times 720$ h), cifra verificable con monitoreo. Sin
embargo, el equipo no tiene acceso a los registros del proveedor y
el texto no define el punto de medición ni qué cuenta como
``mantenimiento programado''. &
Definir el punto de medición (por ejemplo, respuesta exitosa de la
página principal a intervalos fijos), la ventana y el aviso de
mantenimiento y la fuente de datos. Dirigir el requisito al
proveedor del servicio e indicar el origen de la necesidad. \\
\hline

RNF-03 & Requiere corrección &
La meta es medible (90\,\% de participantes, 60 s) y existe un
procedimiento estándar de prueba de usabilidad. Faltan el número
y el perfil de los participantes, el producto a localizar y el
punto de partida. No se vincula por identificador con RF-01 o
RF-02 ni con una necesidad. &
Completar: ``Al menos el 90\,\% de los participantes de la prueba
(número y perfil definidos) localizará, desde la página principal,
un producto asignado en un máximo de 60 s, por búsqueda o por
categorías.'' Vincular con RF-01 y RF-02. \\
\hline

RNF-04 & Requiere corrección &
WCAG 2.2 nivel AA \cite{wcag22} es un criterio externo y auditable
(por eso cumple Ambigüedad y Verificabilidad). Sin embargo, agrupa
numerosos criterios de éxito en un solo enunciado (Singularidad
parcial), ``páginas evaluadas'' no define la muestra
(Completitud) y el cumplimiento total puede depender de componentes
ajenos al control del equipo (Factibilidad parcial). No se vincula a
una necesidad registrada. &
Acotar a un conjunto de páginas (inicio, resultados de búsqueda,
categoría, ficha de producto y carrito) y definir el procedimiento
de auditoría. Si se requiere seguimiento individual, separar por
principio de accesibilidad. \\
\hline

RES-01 & Cumple con ajustes &
Es comprobable: el sitio se abre en un navegador (la barra de
direcciones es visible en las capturas,
Figs.~1 a~9). ``Navegadores'' no indica
cuáles ni qué versiones. &
Reformular: ``La plataforma deberá poder utilizarse desde las
versiones vigentes de los navegadores definidos por el equipo (por
ejemplo, Chrome, Firefox, Edge y Safari).'' \\
\hline

RES-02 & Requiere corrección &
``Normativa ecuatoriana aplicable'' no identifica la norma. La Ley
Orgánica de Protección de Datos Personales \cite{lopdp} y su
Reglamento General \cite{reglamento_lopdp} permiten concretarla,
pero el requisito no enumera obligaciones comprobables, no se
vincula a una necesidad ni a un requisito funcional concreto, y su
estado es ``Restricción propuesta''. &
Citar la ley y el reglamento y enumerar obligaciones verificables
(por ejemplo, la obtención del consentimiento del titular, que el
reglamento regula en su art.~5). Indicar qué requisitos funcionales
tratan datos personales (cuenta de usuario, sesión). \\
\hline

RES-03 & Cumple con ajustes &
Coincide con la sección ``Alcance de la observación'' y se
comprueba porque las evidencias no incluyen pagos reales ni
módulos administrativos. Agrupa tres límites (solo funcionalidades
públicas, sin pagos reales, sin módulos administrativos), y limita
el estudio más que al sistema. &
Separar en tres restricciones o reformular como una sola condición
de alcance. Aclarar que es una restricción del estudio y no del
sistema. \\
\hline

\multicolumn{4}{@{}L{0.9\textwidth}@{}}{\footnotesize\textit{Fuente: Elaboración propia a partir de la matriz de identificación de requisitos, las evidencias de la exploración y las referencias citadas.}}\\
\end{longtable}
\endgroup

\subsection{Hallazgos sobre el conjunto de requisitos}

Además de la evaluación individual, la revisión del conjunto
permitió identificar los siguientes aspectos:

\begin{itemize}
    \item \textbf{Necesidad sin requisito.} La necesidad
    ``Identificar promociones disponibles'' no tiene ningún
    requisito asociado (RF-01 a RF-08). El sitio muestra una
    sección de ofertas en su barra de navegación, por lo que
    se recomienda agregar un requisito funcional.
    \item \textbf{Funcionalidades observadas y no cubiertas.}
    Las capturas muestran filtros y ordenamiento en las categorías,
    el enlace para cambiar de tienda y las opciones de entrega
    (entrega a domicilio, retiro en tienda y envío express) con la
    forma de pago. Las dos primeras pueden incorporarse a RF-02 y
    RF-08. Las opciones de entrega y pago quedan fuera del alcance
    declarado (RES-03) y deben registrarse como exclusión.
    \item \textbf{Requisitos no funcionales sin origen.} Las nueve
    necesidades identificadas son de interacción; ninguna
    corresponde a rendimiento, disponibilidad, usabilidad o
    accesibilidad. RNF-01 a RNF-04 deben justificarse como
    propuestas del equipo o vincularse con una necesidad.
    \item \textbf{Prioridad.} Ningún requisito indica su
    prioridad, lo que dificulta decidir el orden de corrección.
\end{itemize}

\subsection{Actores del proceso}

El proceso de requisitos del sistema involucra cuatro actores.
Para cada uno se distingue si participa de forma real en este
estudio o si se propone como rol necesario para continuar el
proceso.

\begingroup
\footnotesize
\setlength{\tabcolsep}{2.5pt}
\renewcommand{\arraystretch}{1.15}

\begin{longtable}{|L{0.12\textwidth}|L{0.21\textwidth}|L{0.27\textwidth}|L{0.27\textwidth}|}
\caption{Actores del proceso de requisitos}
\label{tab:actores}\\
\hline
\textbf{Actor} & \textbf{Condición} & \textbf{Función en el proceso} &
\textbf{Relación con el sistema} \\
\hline
\endfirsthead

\multicolumn{4}{c}{\textit{Continuación de la Tabla~\ref{tab:actores}}}\\
\hline
\textbf{Actor} & \textbf{Condición} & \textbf{Función en el proceso} &
\textbf{Relación con el sistema} \\
\hline
\endhead

Ingeniero de requisitos &
Real: equipo de estudiantes que realiza la práctica. &
Elicita las necesidades a partir de la exploración, las analiza,
las especifica en la matriz de requisitos y las valida con las
propiedades de calidad. &
Autor de la matriz de identificación de requisitos y de la
evaluación de calidad. Trabaja sobre el sistema publicado, sin
acceso a su código ni a sus módulos internos. \\
\hline

Cliente &
Real: Almacenes Tía como organización propietaria de la plataforma.
No se identifica su participación en la exploración. &
Aprobaría los requisitos, definiría prioridades y aceptaría los
resultados de las verificaciones. &
Es dueño del negocio de comercio electrónico; los requisitos
RNF-01 a RNF-04 y RES-02 le corresponden como objetivos o
condiciones de negocio. \\
\hline

Usuario &
Real: consumidor que usa la plataforma. Las capturas muestran una
sesión de usuario autenticada. Propuesto: participantes de la
prueba de usabilidad de RNF-03. &
Fuente de las necesidades de búsqueda, navegación, carrito y
retiro; valida que los requisitos describen su uso real. &
Usa directamente RF-01 a RF-08. En el contexto local, son
consumidores de San Carlos, cantón Quevedo, que consultan
productos y organizan su compra. \\
\hline

Desarrollador &
Real: equipo técnico que mantiene la plataforma; no fue contactado.
Rol propuesto: responsable de implementar los ajustes. &
Evalúa la factibilidad de los requisitos, los implementa y diseña
las pruebas de aceptación. &
Responsable de RF-01 a RF-08 y de medir o cumplir RNF-01, RNF-02,
RNF-04, RES-01 y RES-02. \\
\hline

\multicolumn{4}{@{}L{0.9\textwidth}@{}}{\footnotesize\textit{Fuente: Elaboración propia a partir de \cite{iso29148,pohl2010,swebok4} y del alcance de la exploración.}}\\
\end{longtable}
\endgroup

\subsection{Roles y responsabilidades}

\begingroup
\footnotesize
\setlength{\tabcolsep}{2.5pt}
\renewcommand{\arraystretch}{1.15}

\begin{longtable}{|L{0.13\textwidth}|L{0.5\textwidth}|L{0.28\textwidth}|}
\caption{Responsabilidades de los actores y requisitos relacionados}
\label{tab:roles}\\
\hline
\textbf{Actor} & \textbf{Responsabilidades} & \textbf{Requisitos relacionados} \\
\hline
\endfirsthead

\multicolumn{3}{c}{\textit{Continuación de la Tabla~\ref{tab:roles}}}\\
\hline
\textbf{Actor} & \textbf{Responsabilidades} & \textbf{Requisitos relacionados} \\
\hline
\endhead

Ingeniero de requisitos &
Identificar y clasificar los requisitos; evaluar su calidad;
documentar justificaciones y evidencias; proponer correcciones;
mantener la trazabilidad entre necesidades, requisitos y
evidencias. &
Los quince requisitos (RF-01 a RF-08, RNF-01 a RNF-04 y
RES-01 a RES-03). \\
\hline

Cliente &
Confirmar que los requisitos reflejan sus objetivos de negocio;
definir prioridades; aprobar metas de rendimiento, disponibilidad
y protección de datos; aceptar los resultados. &
RNF-01, RNF-02, RNF-03, RNF-04 y RES-02. \\
\hline

Usuario &
Expresar sus necesidades de uso; participar en pruebas de
usabilidad; validar que las funciones de búsqueda, carrito y
tienda responden a su forma de comprar. &
RF-01 a RF-08 y RNF-03. \\
\hline

Desarrollador &
Estimar la factibilidad; implementar las correcciones; definir
los casos de prueba y los procedimientos de medición;
informar el resultado de cada verificación. &
RF-01 a RF-08, RNF-01, RNF-02, RNF-04, RES-01 y RES-02. \\
\hline

\multicolumn{3}{@{}L{0.9\textwidth}@{}}{\footnotesize\textit{Fuente: Elaboración propia.}}\\
\end{longtable}
\endgroup

\subsubsection{Actores reales y roles propuestos}

En este estudio, los \textbf{actores reales} son el equipo que
realiza la práctica (que actúa como ingeniero de requisitos), el
usuario cuyo uso del sitio quedó registrado en las capturas, y
Almacenes Tía, que existe como organización propietaria y como
equipo de desarrollo pero no intervino en la exploración. Por
esa razón, las necesidades se obtuvieron por observación de la
interfaz y no mediante entrevistas.

Los \textbf{roles propuestos} son los que el proceso requiere para
validar los requisitos y que el estudio no pudo cumplir:
representantes de Almacenes Tía que aprueben requisitos y
prioridades (cliente), participantes de la prueba de usabilidad
de RNF-03 (usuario), un equipo que mida RNF-01, RNF-02 y RNF-04
(desarrollador) y, por RES-02, una persona responsable de la
protección de datos personales. Mientras estos roles no
intervengan, los requisitos no funcionales y las restricciones
normativas continúan como propuestas sin validar.

\subsection{Evidencias}

Las evidencias de esta sección son:

\begin{itemize}
    \item \textbf{Matriz de evaluación de calidad:}
    Tabla~\ref{tab:matriz_calidad}.
    \item \textbf{Justificaciones de las propiedades y correcciones:}
    Tabla~\ref{tab:justificaciones}.
    \item \textbf{Actores y responsabilidades:}
    Tablas~\ref{tab:actores} y~\ref{tab:roles}.
    \item \textbf{Capturas de la exploración} usadas para verificar
    la información de los requisitos funcionales:
    Figs.~1 a~9, en particular la ficha de
    producto con selector de cantidad y confirmación de agregado
    (Figs.~5 y~6) y la categoría Supermercado
    con filtros (Fig.~4).
    \item \textbf{Fuentes utilizadas:} las que se listan a continuación.
\end{itemize}

\subsection{Fuentes utilizadas}

\begin{thebibliography}{9}

\bibitem{iso29148}
ISO/IEC/IEEE, \emph{Systems and software engineering --- Life cycle
processes --- Requirements engineering}, ISO/IEC/IEEE 29148:2018,
Geneva, Switzerland, 2018.

\bibitem{pohl2010}
K. Pohl, \emph{Requirements Engineering: Fundamentals, Principles,
and Techniques}. Berlin, Germany: Springer, 2010.

\bibitem{swebok4}
H. Washizaki, Ed., \emph{Guide to the Software Engineering Body of
Knowledge (SWEBOK Guide)}, version 4.0. Los Alamitos, CA, USA:
IEEE Computer Society, 2024. [Online]. Available:
\url{https://www.computer.org/swebok}

\bibitem{lopdp}
Asamblea Nacional del Ecuador, ``Ley Orgánica de Protección de
Datos Personales,'' \emph{Registro Oficial}, Quinto Suplemento
No.~459, May 26, 2021.

\bibitem{reglamento_lopdp}
Presidencia de la República del Ecuador, ``Reglamento General de la
Ley Orgánica de Protección de Datos Personales,'' Decreto Ejecutivo
No.~904, \emph{Registro Oficial}, Tercer Suplemento No.~435,
Nov. 13, 2023.

\bibitem{wcag22}
W3C, ``Web Content Accessibility Guidelines (WCAG) 2.2,'' W3C
Recommendation, Oct. 5, 2023. [Online]. Available:
\url{https://www.w3.org/TR/WCAG22/}

\end{thebibliography}

\end{document}
